\section{Methodology}

\subsection{Routing Engine}
The routing engine prioritizes safest paths over shortest distance by penalizing high-risk segments derived from crime data.

\begin{figure}[h]
\centering
\includegraphics[width=0.8\textwidth]{"images/SurakṣāMārga.ai landing page.png"}
\caption{Women Safety Navigation Interface}
\end{figure}

\subsection{Risk Scoring Framework}
Crimes targeting women receive maximum analytical weight, guaranteeing algorithmic circumvention of historically hostile streets.

\begin{figure}[h]
\centering
\includegraphics[width=0.8\textwidth]{"images/SurakṣāMārga.ai Map page.png"}
\caption{Safe Route Visualization for Women}
\end{figure}

\subsection{Temporal Intelligence}
The system enforces active multipliers for night-time navigation, redirecting predictive paths heavily toward well-lit, trafficked thoroughfares.

\begin{figure}[h]
\centering
\includegraphics[width=0.8\textwidth]{"images/SurakṣāMārga.ai feature in map.png"}
\caption{Safety-First Feature Mapping}
\end{figure}

\subsection{SOS Emergency System}
An integrated long-press instantly bypasses standard navigation to initiate live geo-tracking and transmit high-priority alert loops to responder networks.

\begin{figure}[h]
\centering
\includegraphics[width=0.8\textwidth]{"images/SurakṣāMārga.ai safe route found.png"}
\caption{Safe Route Output Generation}
\end{figure}

\subsection{5G Safety Layer}
Operating under 5G capabilities drastically removes data transmission limits, facilitating split-second distress beacon broadcasts without buffer wait times.

\begin{figure}[h]
\centering
\includegraphics[width=0.8\textwidth]{"images/SurakṣāMārga.ai SAFETY LOOP CLOSURE Flow:START → Route Selected → Tracking → Arrival → SAFE CONFIRMATION.png"}
\caption{SOS Response and Real-Time Safety Closure Flow}
\end{figure}

\subsection{Data Scope and Deployment Limitation}
\textbf{Note on Data Scope:} For demonstration purposes and due to current data availability constraints, the system presently generates and recommends safe routes only within Bengaluru. The underlying architecture is designed to be scalable, and future work will focus on expanding the system to additional cities as more comprehensive datasets become available.
